\section{Threats to validity}
\label{sec:threats}

\subsection{Construct validity}
\label{sec:threats:construct}

\paragraph{Authorship is not binary.} The strongest threat to this study is that its
outcome variable does not exist in the form the measurement assumes. Real contributions
are frequently neither purely human nor purely machine: agent output is edited by
humans before merge, and human code is refactored by agents. Treating authorship as a
binary label imposes a dichotomy on a continuum. We mitigate this in three ways. Where
possible we use the evidence-typed representation of Section~\ref{sec:provenance}
instead of a binary label; we report a sensitivity analysis excluding heavily edited
pull requests, using AIDev's review metadata; and we frame all claims as concerning
\emph{declared} agent authorship rather than AI involvement in general.

\paragraph{The unobservable stratum.} Our positive set captures agentic contributions
that platforms declare. It cannot capture a developer who pastes generated code from a
chat interface into their own editor, which is plausibly the most common form of AI
assistance in practice. This study measures detector validity on the observable
stratum, and is explicit that the unobservable one is out of scope. The N1--N2
comparison in Section~\ref{sec:groundtruth:n2} offers indirect evidence about its
magnitude under stated assumptions, but this is a bound, not a measurement.

\subsection{Internal validity}
\label{sec:threats:internal}

\paragraph{Negative-set purity.} N1's purity rests on chronology plus the requirement
that artifacts were not modified after the cutoff. Rewritten repository history could in
principle carry post-cutoff text under pre-cutoff timestamps; we report the prevalence
of detectable history rewriting in the sample. The cross-check against archived event
streams described in Section~\ref{sec:groundtruth:n1} tests the timestamp rule directly
rather than assuming it.

\paragraph{Quiescence bias.} Requiring that artifacts were never modified after the
cutoff selects for those that attracted no later activity. Such artifacts may be
systematically shorter, less discussed or from less active projects, which could bias
error-rate estimates in either direction. \RESULT{Report the observable differences
between eligible and ineligible pre-cutoff artifacts on length, project activity and
genre, so readers can judge the direction of the bias.}

\paragraph{Repository eligibility.} Only repositories predating the cutoff support
matched negatives, excluding a substantial minority of the frame
(Section~\ref{sec:groundtruth:n1}). Excluded repositories are newer and may differ in
agent adoption; the generalisation check on positives from excluded repositories tests
whether this matters.

\paragraph{Detector version drift.} Detector implementations and their underlying
scoring models change. All versions are pinned and reported, and the replication package
fixes them. This is also why no commercial detector is included: its behaviour cannot be
pinned, so a result obtained from it cannot be reproduced.

\subsection{External validity}
\label{sec:threats:external}

\paragraph{The seven-month window.} All positives fall between \AidevWindowStart{} and
\AidevWindowEnd{}. Agent platforms and their underlying models change rapidly, and
detector performance against a 2025 model generation may not transfer to later ones.
This bounds the shelf life of the specific error rates while leaving the methodological
contribution --- the counterfactual design and the correction procedure --- applicable
to whatever generation follows.

\paragraph{Agent coverage.} AIDev covers agent platforms with public GitHub footprints.
Agents operating under human accounts, and self-hosted or enterprise deployments, are
absent. Combined with the \AidevTopAgentShare{} share of the largest agent, this makes
per-agent reporting essential rather than optional.

\paragraph{Repository population.} The curated frame consists of comparatively popular
repositories, with a median of \RepoStarsMedian{} stars. Detector behaviour on small,
inactive or personal projects is not established by this study.

\subsection{Conclusion validity}
\label{sec:threats:conclusion}

The stratified analysis involves many comparisons across genres, languages, agents and
length bins. We report effect sizes with confidence intervals throughout and avoid
null-hypothesis significance language, in keeping with the reporting conventions of this
journal. Where a stratum is too thin to support a precise estimate, we report the
interval rather than the point and say so explicitly.
